Time-series classifiers deployed in sensor-driven environments often encounter corrupted inputs such as missing segments, local noise, and amplitude drift. These corruptions can make prediction sets from standard post-hoc uncertainty methods overly large or poorly matched to the observed input condition. This paper presents \fastcp, a lightweight post-hoc conformal prediction framework for corrupted time-series classification under limited compute. The method computes cheap perturbation fingerprints, estimates a fingerprint-weighted local threshold, and mixes it with an augmented global conformal threshold. This design supports two practical operating modes: a coverage-prioritized mode that falls back to global augmented calibration, and a Pareto-oriented efficient mode that accepts a small measured coverage cost in exchange for smaller prediction sets. On 35 small-to-medium UCR datasets with five random seeds and three corruption families, the efficient setting reduces average prediction-set size by 3.02--4.44\% relative to augmented split conformal prediction at $\alpha=0.05$, with a mean coverage decrease of 0.17--0.86 percentage points. A mismatched-corruption stress test and dataset-level block bootstrap further suggest that the set-size reduction is not solely an artifact of matched corruptions or cell-level dependence. Additional UEA multivariate, MiniROCKET, and FCN checks provide supporting evidence that the qualitative tradeoff is not tied to a single backbone. \fastcp{} is therefore intended as a reproducible, deployable coverage-efficiency tradeoff layer, not as a new conformal-validity guarantee under distribution shift.
